% ===== PRÉAMBULE CANONIQUE — à coller VERBATIM en tête de CHAQUE .tex =====
\documentclass[11pt,a4paper]{article}
\usepackage[utf8]{inputenc}\usepackage[T1]{fontenc}\usepackage{lmodern}
\usepackage[french]{babel}
\usepackage{amsmath,amssymb,mathtools}\usepackage{icomma}
\usepackage{siunitx}\sisetup{locale=FR}  % \qty{}{}, \si{}, \ang{}
\usepackage{xcolor}\usepackage{colortbl}
\usepackage{tikz}\usepackage{pgfplots}\pgfplotsset{compat=1.18}
\usepackage[siunitx,europeanresistors]{circuitikz} % circuits (élec.)
\usepackage[most]{tcolorbox}\usepackage{enumitem}\usepackage{array,booktabs}
\usepackage[a4paper,margin=2.2cm]{geometry}
\usetikzlibrary{arrows.meta,calc,angles,quotes,positioning,patterns,%
  backgrounds,shapes.geometric,decorations.pathmorphing,decorations.markings,intersections}
\definecolor{primary}{RGB}{222,49,99}\definecolor{primarydark}{RGB}{150,22,55}
\definecolor{defcol}{RGB}{0,114,178}\definecolor{methcol}{RGB}{52,92,148}
\definecolor{excol}{RGB}{0,135,90}\definecolor{attncol}{RGB}{200,40,55}
\tcbset{boxbase/.style={enhanced,breakable,boxrule=0.9pt,arc=2.5pt,
  left=3mm,right=3mm,top=2.3mm,bottom=2.3mm,fonttitle=\bfseries,coltitle=white}}
% --- boîtes TUTORIEL (noms EXACTS, ne pas renommer) ---
\newtcolorbox{cle}[1][]{boxbase,colback=defcol!6,colframe=defcol,
  title={\textsc{À retenir}\ifx\relax#1\relax\relax\else\textmd{ : \itshape #1}\fi}}
\newtcolorbox{methode}[1][]{boxbase,colback=methcol!7,colframe=methcol,sharp corners,
  title={\textsc{Méthode}\ifx\relax#1\relax\relax\else\textmd{ --- \itshape #1}\fi}}
\newtcolorbox{exemplebox}[1][]{boxbase,colback=excol!5,colframe=excol,
  title={\textsc{Exemple}\ifx\relax#1\relax\relax\else\textmd{ : \itshape #1}\fi}}
\newtcolorbox{piege}[1][]{boxbase,colback=attncol!6,colframe=attncol,
  title={\textsc{Piège}\ifx\relax#1\relax\relax\else\textmd{ : \itshape #1}\fi}}
\newtcolorbox{remarque}[1][]{boxbase,colback=black!4,colframe=black!45,
  title={\textsc{Remarque}\ifx\relax#1\relax\relax\else\textmd{ : \itshape #1}\fi}}
% --- boîtes EXERCICE / CORRIGÉ (noms EXACTS, auto-comptées) ---
\newtcolorbox[auto counter]{exo}[1][]{boxbase,colback=primary!4,colframe=primary,
  title={\textsc{Exercice}~\thetcbcounter\ifx\relax#1\relax\relax\else\textmd{ --- \itshape #1}\fi}}
\newtcolorbox[auto counter]{corrige}[1][]{boxbase,colback=excol!4,colframe=excol,
  title={\textsc{Corrigé}~\thetcbcounter\ifx\relax#1\relax\relax\else\textmd{ --- \itshape #1}\fi}}
\newcommand{\vv}[1]{\vec{#1}}\newcommand{\ds}{\displaystyle}
\newcommand{\R}{\mathbb{R}}\newcommand{\N}{\mathbb{N}}\newcommand{\Z}{\mathbb{Z}}
% =========================================================================
\begin{document}
\sisetup{per-mode=symbol}

\begin{exo}[décomposer le poids sur une pente]
Un bloc de masse $m=\qty{10}{\kilogram}$ est posé sur un plan incliné à $\alpha=\ang{30}$. On prend
$g=\qty{10}{\metre\per\second\squared}$, $\sin\ang{30}=\num{0,5}$, $\cos\ang{30}\approx\num{0,87}$.
\begin{center}
\begin{tikzpicture}[>=Stealth,scale=0.85]
  \coordinate (O) at (0,0); \coordinate (B) at (5,0); \coordinate (C) at (5,2.887);
  \draw[thick,fill=black!5] (O)--(B)--(C)--cycle;
  \draw (0.8,0) arc (0:30:0.8); \node at (1.1,0.2){$\alpha$};
  \coordinate (P) at (2.6,1.5); \fill[defcol] (P) circle (0.12);
  \draw[->,line width=1pt,attncol] (P) -- ++(0,-1.5) node[below]{$\vv G$};
\end{tikzpicture}
\end{center}
\begin{enumerate}[label=\alph*)]
  \item Calcule le poids $G$ du bloc.
  \item Calcule la composante \emph{motrice} $G_x=mg\sin\alpha$ (le long du plan).
  \item Calcule la composante plaquante $G_y=mg\cos\alpha$ (perpendiculaire au plan).
\end{enumerate}
\end{exo}

\begin{exo}[la réaction normale]
On prend $g=\qty{10}{\metre\per\second\squared}$, $\cos\ang{30}\approx\num{0,87}$.
\begin{enumerate}[label=\alph*)]
  \item Un bloc de $\qty{10}{\kilogram}$ repose sur un sol \emph{horizontal}. Que vaut la réaction normale $N$ ?
  \item Le même bloc est posé sur un plan incliné à \ang{30}. Que vaut alors $N=mg\cos\alpha$ ?
\end{enumerate}
\end{exo}

\begin{exo}[le frottement cinétique]
Un bloc de masse $m=\qty{10}{\kilogram}$ glisse sur un sol horizontal. Le coefficient de frottement
cinétique vaut $\mu_c=\num{0,3}$ et $g=\qty{10}{\metre\per\second\squared}$.
\begin{enumerate}[label=\alph*)]
  \item Que vaut la réaction normale $N$ ?
  \item En déduis l'intensité de la force de frottement cinétique $F_f=\mu_c N$.
\end{enumerate}
\end{exo}

\begin{exo}[descente sans frottement]
Un palet est lâché, sans vitesse initiale, sur un plan incliné à $\alpha=\ang{30}$, \emph{sans frottement}
($g=\qty{10}{\metre\per\second\squared}$).
\begin{enumerate}[label=\alph*)]
  \item Montre que son accélération le long de la pente vaut $a=g\sin\alpha$, et calcule-la.
  \item Cette accélération dépend-elle de la masse du palet ?
\end{enumerate}
\end{exo}

\begin{exo}[glissement à vitesse constante]
On tire une caisse sur un sol horizontal, à \emph{vitesse constante}, avec une force horizontale
$\vv F=\qty{40}{\newton}$.
\begin{enumerate}[label=\alph*)]
  \item Que vaut l'accélération de la caisse ?
  \item En déduis l'intensité de la force de frottement.
\end{enumerate}
\end{exo}
\end{document}
